\documentclass[UTF8,a4paper,10pt,fontset=none]{ctexart}
\usepackage[margin=18mm,headheight=15pt]{geometry}
\usepackage{amsmath,amssymb,booktabs,array,tabularx,xcolor,fancyhdr,hyperref,enumitem}
\setmainfont{Arial}
\setCJKmainfont{Microsoft YaHei}
\setCJKmonofont{Microsoft YaHei}
\definecolor{ink}{HTML}{18344C}
\hypersetup{unicode=true,colorlinks=true,urlcolor=ink,linkcolor=ink,
 pdftitle={YunTAPR-Net 极端降水诊断准备与预注册候选方案},
 pdfauthor={YunTAPR-Net Research}}
\pagestyle{fancy}\fancyhf{}
\lhead{\small YunTAPR-Net}\rhead{\small 诊断准备完成／真实新增计算未授权}
\cfoot{\small 2026-10-31\quad|\quad\thepage}
\setlength{\parindent}{0pt}\setlength{\parskip}{4pt}
\renewcommand{\arraystretch}{1.15}\setlength{\tabcolsep}{5pt}
\setlist{nosep,leftmargin=1.8em,topsep=3pt}
\ctexset{section={format=\large\bfseries,beforeskip=9pt,afterskip=4pt},
 subsection={format=\normalsize\bfseries,beforeskip=6pt,afterskip=3pt}}
\emergencystretch=3em
\newcommand{\key}[1]{\textcolor{ink}{\textbf{#1}}}
\newcommand{\code}[1]{{\small\texttt{\detokenize{#1}}}}
\newcommand{\repo}{https://github.com/flyeagle-cex/YunTAPR-Net/blob/7c19b9d36bc73b998fcf72b9cc630989f7e3b324/}
\begin{document}
{\LARGE\bfseries 极端降水诊断准备}\\[4pt]
{\Large 零新增训练的预注册候选方案}\\[6pt]
2026年10月31日\quad 版本V1\quad 研究者主导，AI辅助

\section*{1. 本轮已经做了什么}
\key{先区分强雨点值转换、条件尾部分布与事件／观测证据缺口，再决定后续投入。}
研究者指令“按你说的来吧”“直接开始操作，不需要给我看了”用于推进准备、已有证据诊断、候选工具与合成验证；本文件不伪造为研究者已冻结所有方法或已批准新增推理预算。

基线为工作分支\code{prep/b2-stage-acceptance-20261025}的
\code{7c19b9d36bc73b998fcf72b9cc630989f7e3b324}，获取远端后本地与远端一致。
所有源输入按该提交中的Git对象字节绑定，避免Windows换行差异影响溯源。

\begin{center}\small
\begin{tabularx}{\linewidth}{@{}Xr@{}}
\toprule
本轮针对性检查 & 结果\\
\midrule
总体五指标与源JSON定位指针逐项核对 & 178项一致\\
已有配对差值独立相减、参考值再核对 & 5,737项一致\\
方向性对照使用明确B7／PERP源值重新相减 & 65项一致\\
九模型、A/B/C、严格$>20$分类计数核对 & 27组一致\\
九份固定E9身份与矩阵／已有字节数收据绑定 & 9份一致\\
历史19份补算结果及其原指标一致性门槛 & 19份通过，继承\\
候选概率界的独立合成检查 & 14项通过\\
\bottomrule
\end{tabularx}
\end{center}

\textbf{本轮实际活动：}新增模型推理0、训练更新0、权重读取／反序列化0、原始气象读取0、2025访问0。新工具没有真实数据读取器或实验启动入口；历史模型、损失、协议和收据均保持原样。

\section*{2. 已有事实与不能推断的内容}
已有27组严格$R>20$ mm/h结果均为真阳性TP=0、假阳性FP=0、漏报FN=717、临界成功指数CSI=0、命中率POD=0。无预测阳性时虚警比FAR分母为0，保留未定义。

717个像元时次来自共同2024开发人口，分布于207个采样窗口；$\ge20$为208个窗口，边界不能改写。$>50$仅6个相关像元时次、5个采样窗口，独立天气事件数未知。

保存的最高条件分位数q32最大值换算后约41.43--100.32 mm/h，否定“输出统一硬截断在20”的解释；它们是否位于真实强雨位置仍缺联合证据。候选点估计失败不能提升为全部概率预测失败。

\key{当前状态：准备完成；新增真实诊断未启动；最终科学结论与2025继续冻结。}

\newpage
\section*{3. 候选范围、对照与定义}
\textbf{模型集合：}仅既有B4、B7、B8各2029/2030/2031随机种子，共九份E0、固定第九轮末检查点E9 LAST。恢复仍属于原种子，不挑最佳种子，不集成，不训练控制或新模块。B0/B1的不同人口不进入本轮核心排名；ZERO/PERP已有结果只作背景证据。

\textbf{数据集合：}2023训练959场景、2024开发963场景。开发有效监督3,303,090个像元时次，有雨407,820个。监督使用冻结行政区与数据有效性交集；地下压力层或DOTE（动态地形与热力／水汽强迫表征）特征无效只能屏蔽相应特征，不能删除原监督像元。

\textbf{时间与单位：}标签窗$[T,T+30\mathrm{min})$，分析时刻$T+30\mathrm{min}$，仅00/06/12/18 UTC（协调世界时）；GFS（全球预报系统）f006为六小时有效预报的回顾估计，不证明历史业务可用性。标签单位mm/h；30分钟累计量为雨强乘0.5小时。

\subsection*{拟提交研究者冻结的规则}
\begin{enumerate}
\item 主阈值保留严格$R>20$ mm/h；5、10用于桥接，50为稀少尾部次要诊断。额外保存标签和点预测恰等于20、50的计数，不重命名旧分层。
\item 主比较B8减B7；次要比较B7减B4、B8减B4。每种子先配对，再报告三种子均值、样本标准差和改善／退化次数。种子不是独立天气事件。
\item A/B/C三种点估计全部保留为诊断代理；不按开发表现挑选最终估计。概率决策0.5只作历史参考，不代表灾害业务最优阈值。
\item 主概率诊断使用有限分位数支持的概率区间；不把未知尾部填0，不拟合新尾部分布，不让模型只在各自支持子集上排名。
\item 全部32个条件分位数覆盖、七个原雨强层、四UTC及已有三高程组均保留。缺失联合分层明确标缺失，不从边际结果反推。
\item 2023正雨P95/P99（第95／99百分位）拟以逐像元$R>0$为主口径，$R>0.1$为敏感性口径，使用线性样本分位定义。数值当前未计算；不使用开发分位或逐场景最大值分位替代。
\end{enumerate}

\subsection*{参考基线与阐释范围}
发生概率$p=P(R>0.1\mid X)$，雨强条件分布针对$R>0.1$。
超阈常数气候基线拟用同人口2023训练监督上的$P(R>r)$，所有模型共享；需要一次经授权的训练标签统计，不能使用2024的事件频率拟合基线。

\key{这些是候选方法规则，不修改已有冻结规范，也不代表研究者已批准新计算。}
在完整分布未识别时，诚实交付“不确定／未识别”也是有效诊断结果。

\newpage
\section*{4. 概率候选工具：保留无法识别的部分}
设$z=\log(1+r)$，32个严格有序的条件对数分位数为$q_i$，
$\tau_i=(i-0.5)/32$。仅从分位数语义与分布函数单调性可得：
\[
 L_F(z)=\max\bigl(\{0\}\cup\{\tau_i:q_i\le z\}\bigr),\qquad
 U_F(z)=\min\bigl(\{1\}\cup\{\tau_i:q_i>z\}\bigr).
\]
\[
 \underbrace{p(1-U_F)}_{\ell}\ \le\
 \widehat P(R>r\mid X)=p\{1-F_+(z\mid X)\}
 \ \le\ \underbrace{p(1-L_F)}_{u}.
\]
这些是\textbf{与模型输出相容的分布范围}，不是实际风险置信区间，不保证模型校准。
当$r=0.1$时按冻结支持语义直接返回$[p,p]$。若$z\ge q_{32}$，则
\[
0\le\widehat P(R>r\mid X)\le p/64;
\]
不能因为下界为0就把真实模型尾概率写成0。

\subsection*{完整人口的指标界}
Brier Score（布里尔分数）针对超阈标签$y\in\{0,1\}$。每像元的平方误差界为
\[
 b_{\min}=\begin{cases}\ell^2&y=0\\(1-u)^2&y=1\end{cases},\qquad
 b_{\max}=\begin{cases}u^2&y=0\\(1-\ell)^2&y=1.\end{cases}
\]
对完整监督人口取均值；比较A、B时使用
$[BS_{A,\min}-BS_{B,\max},\,BS_{A,\max}-BS_{B,\min}]$。
区间跨0表示现证据不能确定方向，不能挑一个尾部补全来宣布胜负。

AUROC（受试者工作特征曲线下面积）的候选界使用正负样本两两排序：
下界比较正样本$\ell$与负样本$u$，上界反向，平分计0.5。
工具用排序与二分计数，避免建立数十亿像元对矩阵。
只有单一类别时保留未定义。

AP，Average Precision（平均精确率）在未知完整分数时不报告伪精确值。
支持域线性插值或其他完整分布补全如需评价，必须作为另外冻结的敏感性分析，不能静默进入主指标。

\subsection*{工程验证范围}
候选工具\code{extreme_diagnostic_candidates.py}已通过14项合成检查：
已知对数均匀分布、q32边界、未知尾部、零发生概率、阈值单调性、非法输入、
Brier穷举端点、含并列的AUROC穷举端点、单类别及决策未识别。
\textbf{尚未用于任何真实模型输出；没有新增气象读取或评价。}

\newpage
\section*{5. 联合输出与诊断报告应保存什么}
未来如获准确授权，拟每份既有E9只做一次共同开发人口提取；原冻结评价器继续作为原指标核对依据。每模型原Brier、AUROC、AP、Conditional Pinball Loss（条件分位数损失）、q32及整数计数全部通过后，才接受新增诊断。

\subsection*{必要联合信息}
共同的场景ID、分析时间、像元／网格ID、物理雨强标签、监督有效性、发生概率$p$及32个FP64（64位浮点）对数分位数必须能关联。
标签和坐标可共享保存，不为每模型重复复制。
若只存充分量，也必须足以恢复预声明比较、误差分层及追溯；只存发生概率正负计数无法恢复极端标签或空间事件。

保存阈值支持域、概率区间宽度、各阈值边界相等计数；高程、压力层可用性、迎／背风、水汽联合诊断需要相应同源条件键。不能用少量工程缓存代表963场景。

\subsection*{条件分布与点值分别报告}
q32名义覆盖为$63/64=98.4375\%$，有雨总体经验覆盖与名义差异用于描述校准。
按真实雨强条件化的分层覆盖不具有逐层同一名义目标。
全部32条覆盖、最高分位欠幅与原条件Pinball一起保留；Pinball为对数雨强尺度，不能冒充mm/h误差。

A=$p\,\mathrm{mean}_i[\exp(q_i)-1]$，B为原混合中位数，
C=$p\{\exp[\mathrm{mean}_i(q_i)]-1\}$。
C不是条件算术均值；A也只是有限分位点代理。
三者各报告固定阈值分类、强雨条件MAE（平均绝对误差）、RMSE（均方根误差）、
Bias（预测减观测的平均偏差）和峰值位置误差，不选开发最优者充当最终泛化。

\subsection*{事件与独立观测的边界}
四UTC稀疏场景、207／208个窗口、35周块和相关像元都不等于独立暴雨。
现阶段只报告采样窗口／像元时次；没有高频观测与追踪键时不生成独立事件结论。

真正的事件评价须先明确时间连接、空间连通、系统分裂合并、边界删失、时间／空间容差及独立观测支持尺度。
这些规则与数据来源需要研究者另外决定，当前不自行指定某个小时或距离阈值。
独立雨量站／雷达等观测也需质量、时间窗及空间支持匹配。
新提取九份模型输出本身不能解决事件独立性问题。

\key{完整报告同时保留改善、退化、未定义与未识别；成功标准不要求出现正增益。}

\newpage
\section*{6. 分阶段推进与资源边界}
\begin{description}
\item[G0：本轮已完成] 保存证据核对、九模型身份清单、候选概率工具及合成检查、可修改方案。无真实新评价。
\item[G1：真实读取前] 研究者冻结方法；绑定具体执行适配器、提交、唯一输出目录、读取清单、硬时间／存储上限和准确新授权。当前这些准入尚未闭合。
\item[G2：需新授权] 2023共享标签统计一次；九份E9各一次2024开发提取。旧原指标回归门槛和整数计数全部通过，否则停止相关后续，不自动重放。
\item[G3：诊断交付] 全部种子、对照、原值、差值、分层、识别区间与负结果齐全。未识别范围如实报告，不为补出结论更换尾部模型。
\item[G4：另行决策] 是否取独立观测、开展物理反驳对照或多尺度DTFM（动态地形特征调制）；训练和2025均不在本方案授权内。
\end{description}

\subsection*{已有清单可推导的资源参考}
\begin{center}\small
\begin{tabularx}{\linewidth}{@{}Xrr@{}}
\toprule
逻辑内容类别 & 字节 & 实例\\
\midrule
九次开发天气／掩膜 & 239,753,278,008 & 71,550\\
六次静态包 & 1,876,968 & 6\\
一次2023标签与掩膜 & 180,049,094 & 246\\
\midrule
上述类别合计 & 239,935,204,070 & 71,802\\
\bottomrule
\end{tabularx}
\end{center}
训练标签参考来自245份既有IMERG（多卫星融合降水产品）文件180,028,178字节，另加掩膜20,916字节。
九份检查点归档长度合计482,176,887字节；如各哈希一次、加载一次，归档长度逻辑参考964,353,774字节另列。
这不承诺\code{torch/zip}内部物理读取量；当前未打开这些权重。

历史九次评价总耗时11,690.4秒，仅估时。新增输出和界计算的硬时间、磁盘空间、最低可用内存、适配器身份尚未冻结；不能把历史6小时训练额度或该估时当新推理预算。重试与重放候选均为0。

\subsection*{失败和范围控制}
人口／检查点身份不一致、数值非有限、原指标回归失败、源读取异常、时间／存储超限时保留部分输出和账本，停止相关工作。
不删数据，不覆盖旧结果，不跳过困难样本，不以未用完额度自动恢复。
已取消的自动化不重开，GitHub只推工作分支，不合并main。

\newpage
\section*{7. 三个可证伪问题与后续取舍}
\textbf{H1：点值决策掩盖风险排序。}
已有点值全部漏检，但分位输出可超过20。若冻结的完整人口概率界与参考比较仍不能支持风险优势，就不能宣称已识别强雨风险；若区间太宽，应报告未识别，不拟合有利尾部。
最小证据为同位置$y,p,q$；需新的有限推理授权，不需训练。

\textbf{H2：条件尾部不足与稀有样本／目标有关。}
q32欠覆盖是事实，稀有强雨、对数尺度和优化竞争是候选原因。
若独立匹配评价中条件尾部正常而问题仅发生在点值转换，则该解释受反证。
先看固定分层损失、欠幅、跨度输出与训练正雨分布；诊断前不改损失、不补种子。

\textbf{H3：观测支持尺度／时间代表性影响极端判断。}
当前没有独立事件与观测链。若严格匹配的独立观测上仍出现同类失败、且产品差异不足以解释，则该解释受反证。
需要另行选择观测与事件方法并批准数据／计算；不能用观测不确定性替模型免责。

\section*{8. 追溯缺陷的处置}
保存结果中65条真实B7减PERP配对行的目标ID与参考ID均填PERP，参考值为空。
本轮明确取真实B7与PERP源值独立相减，65项数值全一致；
缺陷记录在审计JSON，\textbf{未改历史文件}。
旧证据索引仍有PERP未训练的表格行；当前使用完成收据并保留旧展示冲突。

\section*{9. 研究者掌握与论文边界}
研究者本人仍需通过：DOTE（动态地形与热力／水汽强迫表征）的单位／掩膜与混杂；
DTFM有界残差及参数匹配；GFS有效时间与实际可用性；
发生概率和条件分布的组合；同人口同种子同损失的配对与反例。
等级分别为会解释、会实现、会质疑与修改，当前均未考核。

本轮不批准论文结论。已有证据只支持限定人口内的混合预测收益；
DOTE数值效用不等于地形抬升机制，单尺度DTFM不等于完整多尺度设计，
开发年诊断不等于最终泛化。最小下一关是把G1的科学口径、执行身份与准确资源闭合，
其后才有真实新增诊断。

\section*{10. 可追溯材料}
\small
本目录\code{proposal.json}保存候选规则与九个完整E9哈希；
\code{saved_evidence_audit.json}保存178／5,802／27项核对及65条追溯问题；
\code{source_manifest.json}保存基线Git对象字节哈希。
公式候选工具及14项合成检查位于仓库新增工具／测试目录。

源证据：
\href{\repo docs/research_readiness/20261031_extreme_precipitation_audit/EXTREME_PRECIPITATION_AUDIT.md}{极端降水专项审计}；
\href{\repo METRICS.csv}{完整保存指标}；
\href{\repo EXPERIMENT_MATRIX.csv}{实验矩阵}；
\href{\repo docs/research_readiness/20261025/evidence/metric_completion_candidate_started_20261030.json}{19份指标补算收据}；
\href{\repo config/science_v2/quantile_head_v2_frozen_v1.json}{冻结分位数输出语义}。
\end{document}
